\chapter{Slicing and hatching}\label{ch:slicing}

Additive manufacturing processes build a part layer by layer. In two-photon lithography (2PP),
for example, a focused laser writes each layer as a set of \emph{contour lines} along the
boundary and parallel \emph{hatch lines} that fill the interior. \NF{} computes both from the
exact B-rep:

\begin{itemize}
\item \apiref{section} cuts a closed shell with a plane and returns closed, oriented loops of
exact or tightly fitted curves -- not polygons;
\item \apiref{critical_heights} and \apiref{is_extrusion_band} tell you where the topology of
the sections changes and where all layers are identical;
\item \apiref{hatch} fills the region bounded by section loops with hatch lines (intersected
exactly with the loop curves) and inward contour lines (offsets of the loops).
\end{itemize}

All three functions are const queries and thread-safe, so layers can be processed in
parallel.

\section{Sections of solids}\label{sec:section}

\begin{lstlisting}
sec = nf.section(shell, plane, options=nf.SectionOptions())
\end{lstlisting}

\code{section} intersects every face of a closed shell with the plane, clips the
intersection curves to the trimmed regions, chains them into closed loops and orients them
by the outward normals of the solid. The shell is oriented first if necessary.

Compared with \code{planar_section} on a single surface (\cref{sec:planar-section}), the
result is \emph{topologically complete}: loops are closed, holes and islands are
distinguished by orientation, and every piece knows the face it came from.

\subsection{Options}

\begin{center}
\begin{tabularx}{\linewidth}{@{}llX@{}}
\toprule
\apiref{SectionOptions} field & Default & Meaning \\
\midrule
\code{join_tolerance} & $10^{-6}$ & chaining tolerance (model units): ends of neighbouring
pieces closer than this are joined. Use the sewing tolerance of the shell. \\
\code{fit_tolerance} & $0$ & accuracy of traced and fitted section curves on general faces;
$0$ means $10^{-9}$ of the model size. \\
\code{s_axis} & $(0,0,0)$ & first in-plane axis $s$ (projected into the plane and
normalised); zero: chosen automatically. The second axis is $t = n\times s$. \\
\bottomrule
\end{tabularx}
\end{center}

\code{SectionOptions()} takes no constructor arguments; set the fields as attributes:

\begin{lstlisting}
opts = nf.SectionOptions()
opts.s_axis = (1, 0, 0)      # plane coordinates (s, t) = (x, y) for planes normal to z
opts.join_tolerance = 1e-5
\end{lstlisting}

\begin{tip}
For layer-by-layer processing always set \code{s_axis}. With the automatic choice the
in-plane axes are only guaranteed to be orthonormal; fixing them makes the plane coordinates of
all layers comparable, so the hatch direction and phase mean the same in every layer.
\end{tip}

\subsection{The result}

\begin{center}
\begin{tabularx}{\linewidth}{@{}lX@{}}
\toprule
Member & Meaning \\
\midrule
\apiref{PlanarSection}\code{.origin} & origin of the plane coordinates (\code{plane.point}) \\
\code{.s_axis}, \code{.t_axis}, \code{.normal} & orthonormal frame; a plane point is
$p = \mathrm{origin} + s\,\mathbf{s} + t\,\mathbf{t}$ \\
\code{.loops} & list of \apiref{SectionLoop} \\
\midrule
\code{SectionLoop.curves} & consecutive pieces (\apiref{SectionCurve}); the end of one is the
start of the next within the join tolerance \\
\code{SectionLoop.area} & signed area in $(s,t)$: $>0$ for outer loops, $<0$ for holes \\
\midrule
\code{SectionCurve.curve} & the piece in plane coordinates $(s, t, 0)$ \\
\code{SectionCurve.curve3d} & the same piece in 3D (same parameter) \\
\code{SectionCurve.pcurve} & the piece as $(u, v, 0)$ on the surface of its face (same
parameter) \\
\code{SectionCurve.face} & index of the source face in the shell \\
\code{SectionCurve.exact} & \code{True} for an exact isoparametric curve, \code{False} for a
traced and fitted curve \\
\code{SectionCurve.normal(t)} & outward unit normal of the solid at parameter $t$ \\
\bottomrule
\end{tabularx}
\end{center}

Loops are oriented with the \emph{solid on the left} in $(s,t)$: outer boundaries run
counter-clockwise, holes clockwise, islands inside holes counter-clockwise again. The sum of
all loop areas is therefore the material area of the layer.

Faces whose height along the plane normal depends on one parameter only -- planar faces,
extrusions along the normal and surfaces of revolution about it -- give \emph{exact}
section curves: an iso-parameter curve of the face. Other faces give curves traced by the
surface/surface intersection and fitted to \code{fit_tolerance}.

\subsection{Example}

The following script builds a plate with a spherical dome and a bore, sections it at two
heights, and investigates its critical heights:

\example{12_sections.py}

\begin{lstlisting}[style=text]
body: 9 faces, volume 31.330383
origin Vec3(0.000000, 0.000000, 0.500000), s Vec3(1.000000, 0.000000, 0.000000), t Vec3(0.000000, 1.000000, 0.000000), normal Vec3(0.000000, 0.000000, 1.000000)
loop 0: area +24.000000 (outer), 4 piece(s): face 2 exact, face 5 exact, face 3 exact, face 4 exact
loop 1: area -0.785398 (hole), 1 piece(s): face 8 exact
material area 23.214602 (24 - pi/4 = 23.214602)
hole piece at t: plane (2.5000, -0.0000), 3D (2.5000, -0.0000, 0.5000), outward normal (-1.000, +0.000, +0.000)
z = 2: 1 loop, area 3.92699082 (pi r^2 = 3.92699082)
critical heights: [-0.5, 0.0, 1.0, 2.5]
0.1..0.9 is an extrusion band: True; 1.1..2.0: False
section at z = 1: GeometryError: section: the plane contains a face (choose a height between critical heights)
\end{lstlisting}

At $z=0.5$ the section consists of the rectangle of the plate (four exact pieces from the four
side faces) and the bore as a clockwise hole. The outward normal of the solid on the bore
points \emph{into} the bore, towards its axis -- away from the material. At $z=2$ only the
dome is cut, giving a circle of radius $\sqrt{1.5^2-1}$.

\subsection{Errors}

\code{section} raises \code{GeometryError} if a chain of pieces does not close (open shell, or
a gap larger than \code{join_tolerance}; the message gives the position) and if the plane
\emph{contains} a face -- planes through critical heights at which a planar face lies in the
section plane are not supported. The example shows the latter at $z=1$, the top of the plate.

\section{Critical heights and extrusion bands}\label{sec:critical-heights}

\begin{lstlisting}
hs = nf.critical_heights(shell, direction, tolerance=nf.Tolerance())
ok = nf.is_extrusion_band(shell, direction, h0, h1, tolerance=nf.Tolerance())
\end{lstlisting}

\code{critical_heights} returns the sorted heights $h = p\cdot\hat d$ ($\hat d$ = normalised
\code{direction}) at which the topology of the sections can change: vertices, planar faces
normal to the direction, interior extrema of the height on faces (where
$S_u\cdot\hat d = S_v\cdot\hat d = 0$) and on edges ($C'\cdot\hat d = 0$). Between two
neighbouring critical heights the sections have the same number of loops and pieces. Heights
closer than \code{tolerance.length} are merged.

The list is conservative: it may contain heights at which nothing changes. In the example,
$0$ and $1$ are the bottom and top of the plate and $2.5$ is the top pole of the dome. The
sphere (centre $z=1$, radius 1.5) also protrudes below the plate, so $-0.5$ is its lower pole,
the lowest point of the body.

Use the critical heights to
\begin{itemize}
\item avoid section planes \emph{at} a critical height (a plane containing a face raises an
error; a plane through a vertex or touching a face gives degenerate loops) -- slice at
mid-layer heights, or shift by a small amount;
\item adapt the layer thickness, e.g. to place a layer boundary exactly at the top of a
plateau.
\end{itemize}

\code{is_extrusion_band(shell, direction, h0, h1)} returns \code{True} if every face that
meets the open band $h_0 < p\cdot\hat d < h_1$ has a normal perpendicular to the direction
there (within \code{tolerance.angular}). Then the band is an extrusion along the direction and
all its sections are identical: one section can be computed and reused for every layer in the
band. The plate between $z=0.1$ and $0.9$ is such a band; between $1.1$ and $2.0$ the dome is
not.

\section{Hatching}\label{sec:hatch}

\begin{lstlisting}
segments, contours = nf.hatch(loops, direction, spacing, *, phase=0.0, contours=[],
                              min_length=0.0, tolerance=0.0, contour_tolerance=0.0)
\end{lstlisting}

\subsection{Input loops}

\code{loops} may be
\begin{itemize}
\item a \code{PlanarSection} (its loops are used in plane coordinates),
\item a list of \code{SectionLoop}, or
\item a list of closed chains of \code{NurbsCurve} in the $xy$ plane (a single closed curve
counts as a loop); $z$ is ignored.
\end{itemize}

Loops must be oriented with the region on the left: outer loops counter-clockwise, holes
clockwise, islands inside holes counter-clockwise. The region is decided by this
orientation (winding number), \emph{not} by even--odd counting -- for hand-made loops make sure
the holes run clockwise (a circle made with normal $-\hat z$ does).

\subsection{Parameters}

\begin{center}
\begin{tabularx}{\linewidth}{@{}llX@{}}
\toprule
Parameter & Default & Meaning \\
\midrule
\code{direction} & -- & direction $d$ of the hatch lines (2 or 3 components; $z$ ignored,
normalised) \\
\code{spacing} & -- & distance between neighbouring lines ($>0$) \\
\code{phase} & $0$ & offset of the line family: the lines are $n\cdot p =
\mathrm{phase} + k\cdot\mathrm{spacing}$ for integer $k$, with $n = (-d_y, d_x)$, measured from
the origin of the loop coordinates. The same phase gives the same lines in every layer. \\
\code{contours} & \code{[]} & inward offset distances ($>0$) of contour lines; one or more
contours per loop and distance \\
\code{min_length} & $0$ & hatch segments shorter than this are dropped \\
\code{tolerance} & $0$ & contact tolerance: where a loop comes within this distance of a
line without crossing it, the line is not split. $0$: $10^{-9}$ of the model size. \\
\code{contour_tolerance} & $0$ & fit tolerance of the contour curves (see
\code{offset_curve}); $0$: $10^{-7}$ of the model size \\
\bottomrule
\end{tabularx}
\end{center}

\code{hatch} raises \code{GeometryError} if \code{spacing} $\le 0$, the direction is zero in
the $xy$ plane, or more than $10^7$ lines would be needed.

\subsection{Results}

\begin{center}
\begin{tabularx}{\linewidth}{@{}lX@{}}
\toprule
Member & Meaning \\
\midrule
\apiref{HatchSegment}\code{.start}, \code{.end} & end points $(x, y, 0)$; start before end
along the direction \\
\code{.line} & index $k$ of the hatch line \\
\code{.region} & index of the outer loop (area $>0$) of the region containing the segment \\
\code{.start_loop}, \code{.end_loop} & loops on which the segment starts and ends \\
\code{.length} & length of the segment \\
\midrule
\apiref{HatchContour}\code{.curves} & consecutive curve pieces of the closed contour (offset
pieces and arcs about reflex corners) \\
\code{.loop} & first source loop (= \code{loops[0]}) \\
\code{.loops} & sorted indices of all loops the contour was offset from \\
\code{.index} & index of its distance in the \code{contours} list \\
\code{.distance} & the offset distance \\
\bottomrule
\end{tabularx}
\end{center}

Segments are sorted by line index and then along the direction; contours by (first) source
loop and then by the order of the distances.

\subsection{How the lines are computed}

Every hatch line is intersected \emph{exactly} with the loop curves: the signed height
$n\cdot c(t)$ of each Bézier piece is a rational function with the same weights as the
curve, so its sign changes are found by subdividing the control values and refined to
machine precision. There is no polygonisation. Special cases:

\begin{itemize}
\item a loop that only touches a line (within \code{tolerance}) does not split it, so tangent
lines give no zero-length segments;
\item an edge lying exactly on a hatch line is not hatched itself (the line runs along the
boundary). Write such edges as contours.
\end{itemize}

\subsection{How the contours are computed}

A contour at distance $\delta$ is the boundary of the region shrunk by $\delta$: the set of
points of the region at distance $\ge\delta$ from all loops. The raw offset of every loop
(\code{offset_curve} of its pieces into the region, circular arcs about reflex corners) is
split at all self- and mutual intersections, and the fragments that lie inside the region at
distance $\ge\delta$ are chained into closed contours. Consequently
\begin{itemize}
\item contours are trimmed at convex corners;
\item a loop gives \emph{several} contours where the region has a neck narrower than
$2\delta$;
\item the contours of a hole and of another loop \emph{merge} where they come closer than
$2\delta$;
\item a loop gives \emph{no} contour where the region is thinner than $2\delta$ everywhere.
\end{itemize}

\begin{note}[Kinked input curves]
A loop may consist of curves with corners at interior knots, such as a closed polyline or a
\code{make_rectangle_curve} outline (one degree-1 curve). The contours are built from the smooth
pieces between the corners, with arcs at reflex corners, exactly as for a loop made of separate
segments. (Up to version 0.9.0 the offsets of such curves did not close because of an error in
\code{closest_point} at kinks; there the curves had to be split with
\code{curve.decompose_bezier()} first.)
\end{note}

\subsection{Example: hatching a hand-made region}

\example{12_hatch_basics.py}

\begin{lstlisting}[style=text]
10 segments, total length 25.267949
  line 1: (0.0000, 0.2500) -> (4.0000, 0.2500), loops 0->0, region 0
  line 2: (0.0000, 0.5000) -> (4.0000, 0.5000), loops 0->0, region 0
  line 3: (0.0000, 0.7500) -> (0.5670, 0.7500), loops 0->1, region 0
  line 3: (1.4330, 0.7500) -> (4.0000, 0.7500), loops 1->0, region 0
  line 4: (0.0000, 1.0000) -> (0.5000, 1.0000), loops 0->1, region 0
  line 4: (1.5000, 1.0000) -> (4.0000, 1.0000), loops 1->0, region 0
phase 0.1: first line at y = 0.1000, 12 segments
min_length 0.6: 10 of 12 segments kept
45 deg: 20 segments, spacing x length = 7.3304 (region area 7.2146)
contour d = 0.1: from loops [0], 4 pieces, length 11.200000
contour d = 0.3: from loops [0, 1], 5 pieces, length 7.730101
contour d = 0.3: from loops [0, 1], 3 pieces, length 1.073464
contour d = 0.3: from loops [0, 1], 3 pieces, length 1.073464
contour d = 0.1: from loops [1], 4 pieces, length 3.769911
dumbbell, neck width 0.4, d = 0.1: 1 contour(s)
dumbbell, neck width 0.4, d = 0.3: 2 contour(s)
\end{lstlisting}

The output illustrates the rules above:

\begin{itemize}
\item With phase 0 and spacing $0.25$ the lines are $y = 0.25k$. The lines $y=0$ and $y=2$
lie on the edges of the rectangle and are not hatched; lines 3 to 5 are split by the hole.
\item \code{phase=0.1} shifts the family to $y = 0.1 + 0.25k$, giving 8 lines and 12 segments.
\item \code{min_length=0.6} drops the two short pieces left of the hole.
\item The product spacing $\times$ total length approximates the area of the region; it
converges as the spacing decreases.
\item At $\delta=0.1$ each loop has its own contour (the hole's contour has length
$2\pi\cdot0.6$). At $\delta=0.3$ the hole is closer than $2\delta$ to the outer boundary, so
the contours merge, and the narrow gaps split them into three closed contours.
\item The dumbbell has a neck of width $0.4$: one contour at $\delta = 0.1$, two at
$\delta = 0.3 > 0.2$.
\end{itemize}

\section{Workflow: writing paths for two-photon lithography}\label{sec:2pp}

The complete workflow combines the pieces of this chapter: build the part as a closed shell,
section it at the mid-heights of all layers, hatch each layer with contours and alternating
directions, and write the paths in the format the writing system expects.

\example{12_writing_paths.py}

\begin{lstlisting}[style=text]
lens: aperture 40.792 um, sag 8.0 um, volume 5495.6927 um^3, critical heights [0.0, 8.0]
layer  0 z = 0.25: area 1272.149 um^2, 160 hatch lines (5086.03 um), 1 contour(s)
layer  5 z = 2.75: area  903.012 um^2, 136 hatch lines (3613.13 um), 1 contour(s)
layer 10 z = 5.25: area  494.604 um^2, 100 hatch lines (1978.26 um), 1 contour(s)
layer 15 z = 7.75: area   46.928 um^2,  30 hatch lines ( 186.72 um), 1 contour(s)
16 layers, 1794 hatch lines, 16 contours -> lens_paths.txt
hatched volume 5495.89 um^3 vs. cap volume 5495.69 um^3
\end{lstlisting}

Design decisions in the script, and how to adapt them:

\begin{description}
\item[Units.] \NF{} has no unit system: the model is built in micrometres, and every length
(layer distance, spacing, inset) is in micrometres too. Declare the unit when exporting STEP
(\code{unit="um"}, \cref{ch:step}).
\item[Layer heights.] Layers are sectioned at their mid-heights $z = (i+\tfrac12)\,\Delta z$.
The critical heights of the lens are $0$ and $8$ (base and apex); mid-layer planes never hit
them. Check \code{critical_heights} for your part and shift layers that would.
\item[Plane coordinates.] \code{s_axis = (1, 0, 0)} makes the section coordinates equal to
$(x, y)$ for horizontal planes, so the paths can be written directly with the layer height as
$z$. For tilted slicing, map plane coordinates back with
$p = \mathrm{origin} + s\,\mathbf{s} + t\,\mathbf{t}$ or use \code{curve3d}.
\item[Hatch pattern.] The direction alternates between $0^\circ$ and $90^\circ$; the phase of
half a spacing keeps the lines away from the symmetry axes. Other strategies (rotating by
$67^\circ$ per layer, island hatching) only change \code{direction} and \code{phase}.
\item[Contours.] One contour at an inset of $0.15\,\mu$m is written before the hatch lines.
Several distances (\code{contours=[0.1, 0.3]}) give several shells of contour lines.
\item[Meander.] Every second segment is reversed so that the stage does not have to return
to the same side after each line. When a line has several segments (holes), a real path
planner should order segments by region and proximity; the fields \code{line}, \code{region},
\code{start_loop} and \code{end_loop} contain the necessary information.
\item[Export.] Contours are sampled into polylines with \code{NurbsCurve.sample}; hatch
segments are written by their end points. The file format here (one polyline per line:
layer, kind, $x\,y\,z$ triples) is a placeholder -- write whatever your machine reads, e.g.
GWL for Nanoscribe systems. The written volume (spacing $\times$ layer distance $\times$
hatch length) is a useful sanity check: here it matches the cap volume to 0.004\,\%.
\end{description}

\begin{tip}[Performance]
\code{section} and \code{hatch} are const queries, so layers are independent. The Python
binding of \code{hatch} releases the global interpreter lock while it computes, so hatching
many layers in a \code{concurrent.futures.ThreadPoolExecutor} uses several cores;
\code{section} currently keeps the lock, so to parallelise sectioning use a
\code{ProcessPoolExecutor} (shells are rebuilt cheaply in each worker, or sent as STEP text
with \code{to_step}/\code{from_step}). For extrusion bands (\code{is_extrusion_band}) compute
one section and reuse it for all layers in the band.
\end{tip}
